% REPORT::LSPV-3::POINTWISE-DIFFERENTIAL64::2026-10-23T19:24+09:00
\documentclass{article}
\usepackage{amsmath,amssymb,amsthm}
\newtheorem{theorem}{Theorem}
\title{POINTWISE-DIFFERENTIAL64: Arbitrary-Algorithm Bit-Read Coverage}
\author{LSPV-3}
\date{23 October 2026}
\begin{document}\maketitle
\noindent\textbf{Status: P versus NP OPEN.}
\section{Coverage and source audit}
Role-4 PALETTE56 supplies the actual proper-3CNF SAT/UNSAT family; every exact deterministic addressed-bit SAT decider must read q sensitive target positions on the canonical orientation sources, with matching direct linear scan. Role-5 DEGREE96 supplies an identity-inclusion graph-color query bound, also input-order only. Role-1 affine normalization and Role-2 factor joining prove restricted algorithms, not mandatory extracted objects. MAIN85 static context widths do not apply to a one-bit decider. SOURCEPROBE19 is a detached oracle query bound, and exponential regular-resolution and full-DNNF bounds cannot automatically be extracted from arbitrary decider traces. SUPPORT63 proved a coarse global-support simulation; the new theorem below is a strict pointwise refinement. Extraction of the actual read set costs instrumented execution and O(T log N) transcript bits; evaluating derivative incidence externally needs q+1 target encodings. All reads, reduction work, memory, and output are charged.
\section{Three falsifiable constructions}
\begin{description}
\item[Local differential incidence: PROVED/EXACT.] Replace global dependency support by actual bit differences under single source flips.
\item[Balanced full-support encoder: PROVED/EXACT.] Globally nonlocal target orientations nevertheless require q distinct input reads.
\item[Collapsed parity encoder: PROVED/EXACT.] Compute parity in the reducer and leave one informative orientation; target-bit query requirement falls to one.
\end{description}
\section{Pointwise theorem}
Let f(b)=P(e(b)) for a target-bit encoding e. For fixed b*, let I be the set of source coordinates whose individual flips change f. For each target bit j let
\[
J_j=\{i\in I:e(b^*)_j\ne e(b^*\oplus e_i)_j\}.
\]
Define tau as the minimum number of J_j sets covering I.
\begin{theorem}
Every deterministic exact target-bit decider reads at least tau distinct target input bits when executed on e(b*), regardless of its internal algorithm.
\end{theorem}
\begin{proof}
If a sensitive source flip i changes none of the target bits read on e(b*), the target decider has the same adaptive deterministic transcript on e(b* XOR e_i) but the correct answer changes. Contradiction. Thus the read set covers every sensitive source coordinate.
\end{proof}
\section{Proper-3CNF factor-q separation}
For even q at least 4, define H(b)=1 if b=0^q, otherwise zero. Let the balanced orientation be sigma_i=b_i XOR H(b), and the collapsed orientation be sigma_1=XOR_i b_i with all other sigma_i zero. Both preserve parity. Every balanced sigma_i globally depends on all q source bits.
Compose both maps with the original Role-4 PALETTE56 exact proper-3CNF construction of length
\[
n=24q+27,\quad m=k=38q+39,\quad
N=(152q+158)\lceil\log_2(64(n+m+1))\rceil.
\]
Its SAT label is even parity; each one-orientation flip changes exactly six distinct literal words and the changed-word addresses for different orientations are disjoint.
\begin{theorem}
On both b+=1^q (SAT) and b-=1^{q-1}0 (UNSAT), every source flip changes the label. The balanced encoding has tau=q, whereas the collapsed encoding has tau=1. Both are uniform, polynomial-time, all-input label-preserving maps to equal-size valid proper-3CNFs.
\end{theorem}
\begin{proof}
H is zero on the two bases and their one-bit neighbors. Thus a balanced source flip affects its own orientation only, forcing a disjoint target-bit read. A collapsed source flip changes the same first orientation for every input coordinate, so one target-bit position distinguishes each of the q neighboring oppositely labelled inputs.
\end{proof}
\section{Resource ledger, controls, and missing arrow}
Either reducer must read all q source bits to compute its output on a source where parity is sensitive. Balanced routing buffers q bits for one pass or uses two passes and logarithmic input-state space. Collapsed routing requires only a parity accumulator. Both write Theta(N) bits, including fresh padded auxiliary variables, with zero SAT oracle calls and zero recursion. A source-aware canonical recognizer decides the target family in O(q) indexed word operations after full validation. No superlinear decision lower bound follows.
The standard-library reproduction enumerated 336 vectors for q=4,6,8, checked the SAT/UNSAT graph semantics, proper-3CNF witnesses, all one-bit neighbor sources and exact serialized-bit set covers. 2,184 assertion groups passed and two runs gave identical raw JSON, SHA-256 8969860ccaf3aac95dfbb0d6206a107277a28a27e3c6506774788633ee098f93.
The first OPEN arrow remains: an independently defined obstruction must compel superpolynomial charged acquisition/compression/continuation for every uniform exact SAT decider, not only for a source-query or static-compilation model. P versus NP remains OPEN.
\begin{thebibliography}{9}
\bibitem{nisan} N. Nisan, CREW PRAMs and Decision Trees, SIAM J. Comput. 20(6), 999-1007 (1991).
\bibitem{role4} LSPV-4, PALETTE-SOURCE-CYCLE56, repository research/agents/lspv-4 (2026-10-23).
\bibitem{role5} LSPV-5, DEGREE-BOUNDARY96, repository research/main (2026-10-23).
\bibitem{sup63} LSPV-3, SUPPORT-LOCALITY63, repository research/agents/lspv-3 (2026-10-23).
\end{thebibliography}
\end{document}